\documentclass[10pt,a4paper]{amsart}
\usepackage[margin=19mm]{geometry}
\usepackage{amsmath,amssymb,mathtools}
\usepackage[scr=rsfs,cal=euler]{mathalfa}
\usepackage{xfrac}
\usepackage[unicode,hidelinks,bookmarks=false]{hyperref}
\newcommand{\ie}{i.e.\ }
\newcommand{\cf}{cf.\ }
\newcommand{\resp}{respectively\ }
\newcommand{\Subsection}{\subsection}
\providecommand{\nobreakdash}{\nobreak\hbox{-}}
\setlength{\emergencystretch}{2em}
\multlinegap0pt
\usepackage{amssymb}
\usepackage{bm}
\usepackage[shortlabels]{enumitem}
\usepackage{xcolor}
\newtheorem{prop}{Proposition}
\newcommand{\uC}{\underline{\mathbb{C}}}
\title{The Morse index of constant curvature $2$-spheres}
\author{Gavin Ball, Jesse Madnick}
\date{Source: arXiv:2606.06634v2 (CC BY 4.0)}
\begin{document}
\maketitle

\noindent\emph{Source and licence.} The Morse index of constant curvature $2$-spheres. Version arXiv:2606.06634v2 (CC BY 4.0), distributed by arXiv under the Creative Commons Attribution 4.0 International licence, http://creativecommons.org/licenses/by/4.0/, as recorded in the arXiv licence metadata for this exact version and shown on the abstract page https://arxiv.org/abs/2606.06634v2. Full licence notice: \texttt{licenses/arxiv-2606.06634-CC-BY-4.0.txt}.\par\smallskip

\noindent\textit{Editorial scope.} Counted unit: the article's prop prop:normlap (source0.tex lines 377--412). The article's complete original proof of it is delivered with it (source0.tex lines 415--454, one \texttt{proof} environment of 40 lines). No mathematics is written, repaired or re-derived: the statement and the proof are the article's own lines, in the article's own notation, and no retained line is substituted or re-flowed.  Retained uncounted as the source-local context the proof needs: lines 245--255.  Nothing was omitted from any retained span, so the package carries no omission record.  No retained line contains a citation, so the wrapper carries no bibliography and no bibliography entry.  No retained line contains a figure environment, an image-inclusion command or drawing code, so no image is delivered and the package declares no asset dependency.  Editorial formatting only, in addition: an AMS \texttt{amsart} preamble that restates the article's own declarations and macros used by the retained lines, namely the environments \texttt{prop} on separate counters, together with the standard packages those lines need.  The retained environments print in this document as Proposition 1 in order of appearance, whereas the article prints the same items with the numbering of its own full document.  The complete native source of the article as distributed in the arXiv e-print for this exact version is bundled unchanged as \texttt{sources/202199/source0.tex}.
\begin{equation} \label{eq:so3irrmc}
   \mu = \left[ \begin{array}{ c  c | c  c  c  c }
        0 & -\Omega^t & 0 & 0 & \cdots & 0 \\
        \Omega & \Lambda & - c_2 \Psi^t & 0 & \cdots & 0 \\ \hline
        0 & c_2 \Psi & 2 \Lambda & - c_3 \Psi^t & \ddots & 0 \\
        0 & 0 & c_3 \Psi & 3 \Lambda & \ddots & \\
        \vdots & \vdots & 0 & \ddots & \ddots & \\
         & & \vdots & \ddots &(n-1) \Lambda & -c_k \Psi^t \\
         0 & 0 & 0 & & c_k \Psi & n \Lambda
    \end{array}\right]\!,
\end{equation}

\begin{prop}\label{prop:normlap}
    Let $\Delta^\perp$ denote the connection Laplacian associated to the normal connection $\nabla^\perp$ on $N \Sigma_n$, and let $\Delta$ be the connection Laplacian associated to the direct sum connection $\nabla^{(2)} \oplus \cdots \oplus \nabla^{(n)}$ on $N\Sigma_n = \uC_2 \oplus \cdots \oplus \uC_n$.  Then
    \begin{equation*}
        \Delta^{\perp} \begin{bmatrix}
            A_2 \\
            A_3 \\
            \vdots \\
            A_{n-1} \\
            A_{n}
        \end{bmatrix} = \Delta \begin{bmatrix}
            A_2 \\
            A_3 \\
            \vdots \\
            A_{n-1} \\
            A_{n}
        \end{bmatrix} + 4 \, \partial \begin{bmatrix}
            0 \\
            c_3 A_2 \\
            \vdots \\
            c_{n-1} A_{n-2} \\
            c_{n} A_{n-1}
        \end{bmatrix} - 4 \, \overline{\partial} \begin{bmatrix}
            c_3 A_3 \\
            c_4 A_4 \\
            \vdots \\
            c_{n-1} A_{n-1} \\
            0
        \end{bmatrix} - 2 \begin{bmatrix}
            c_3^2 A_2 \\
            (c_3^2 + c_4^2) A_3 \\
            \vdots \\
            (c_{n-1}^2+c_n^2) A_{n-1} \\
            c_n^2 A_n
        \end{bmatrix}.
    \end{equation*}
\end{prop}

\begin{proof}
    From (\ref{eq:so3irrmc}), and switching to complex notation, the connection form for the normal connection is
    \begin{equation*}
        \begin{bmatrix}
           2 i \rho & - c_3 \eta & 0 & & & & \\
            c_3 \overline{\eta} & 3 i \rho & - c_4 \eta & 0 & & & \\
            0 & c_4 \overline{\eta} & 4 i \rho  & -c_5 \eta & \ddots  & &  \\
             & \ddots & \ddots & \ddots & \ddots & & \\
             & & & & (n-1)i \rho & -c_n {\eta} \\
             & & & & c_n \overline{\eta} & n i \rho
        \end{bmatrix}.
    \end{equation*}
    It follows that we have equations
    \begin{equation*}
        \partial^{\perp} = \partial - U, \:\:\:\: \overline{\partial}^\perp = \overline{\partial} + L,
    \end{equation*}
    where $U$ and $L$ are given by
    \begin{equation*}
        U = \begin{bmatrix}
            0 & c_3 & 0 & & \\
            0 & 0 & c_4 & 0 & \\
             & \ddots & \ddots & \ddots & \\
             & & & & c_n \\
             & & & 0 & 0
        \end{bmatrix}, \:\:\:\: L = \begin{bmatrix}
            0 & 0 &  & & \\
            c_3 & 0 & 0 &  & \\
             & c_4 & \ddots & \ddots & \\
             & & \ddots & & \\
             & & & c_n & 0
        \end{bmatrix}.
    \end{equation*} 
We may now calculate
    \begin{align*}
        \Delta^\perp & = 2 \left( \partial^\perp \overline{\partial}^\perp + \overline{\partial}^\perp \partial^\perp \right) \\
        & = 2 \left( (\partial - U)(\overline{\partial} + L) + (\overline{\partial} + L)(\partial - U)  \right) \\
        & = \Delta - 4  \overline{\partial}U + 4  \partial L - 2 \, (UL + LU),
    \end{align*}
where we have used that $U\overline{\partial} = \overline{\partial} U$ and $\partial L = L \partial$.
\end{proof}

\end{document}
